%\vspace{-2.0cm}
\pagenumbering{roman}
\thispagestyle{empty}
\begin{center}
	UNIVERSIDAD CENTRAL DE VENEZUELA\\
	FACULTAD DE CIENCIAS\\
	POSTGRADO EN CIENCIAS DE LA COMPUTACI\'ON\\

	\begin{figure}
						\centering
						  \includegraphics[height=.7\textwidth]{images/UCV.png}
  \end{figure}
  \vspace{1.5cm}
  \large{\textbf{RECUPERACI\'ON, EXTRACCI\'ON Y CLASIFICACI\'ON DE \\ INFORMACI\'ON DE SABER UCV}}

  \vspace{3cm}
  Trabajo de grado de Maestría presentado ante la \\
  ilustre Universidad Central de Venezuela por el\\
  Econ. José Miguel Avendaño Infante para  optar
  al título de \\Magister Scientiarum en Ciencias de la Computaci\'on\\
  \vspace{0.5cm}
  Tutor: Dr. Andres Sanoja\\
  \vspace{1.5cm}
  Caracas - Venezuela\\
  Octubre 2023
\end{center}


\newpage
\thispagestyle{empty}
\large{\textbf{Resumen:}}

Se presenta la investigación \emph{Recuperación, Extracción y Clasificación de Información de Saber UCV} donde se ejecutan procesos de clasificación, almacenamiento y recuperación de información sobre las tesis y trabajos de grado que se encuentran publicados en el repositorio institucional Saber UCV.

Se implementa un Sistema que permite clasificar los documentos mencionados, según el área académica donde cursó estudios el autor de la investigación.

Se extráen los resúmenes de los trabajos y con las clasificaciones obtenidas se conforma un corpus sobre el cual se genera un índice invertido. Al corpus se le aplican técnicas de Procesamiento de Lenguaje Natural, de Minería de Datos y con modelos preentrenados de inteligencia artificial se crean  \textit{embeddings} de textos. Con toda la información procesada se alimenta una base de datos indexada.

El Sistema adicionalmente cuenta con una aplicación web para hacer procesos de Recuperación de Información donde el usuario puede hacer la exploración del corpus, incluida la búsqueda semántica, indicando el: texto a buscar, rango de fechas, área en la cual se generó la investigación, el nivel académico; posteriormente se recuperan los trabajos de mayor relevancia, enriqueciendo la experiencia con la presentación de los resultados en tablas interactivas, Mapas de Conocimiento y recomendaciones de documentos que puedan ser de interés.

La aplicación se implementa bajo un sistema distribuido con la arquitectura cliente-servidor y se soporta en el uso de contenedores orquestados.

\vspace*{2cm}

\textbf{Palabras Clave::} Recuperación de Información, Corpus, Procesamiento del Lenguaje Natural, Relevancia, Inteligencia Artificial, Embeddings, Búsqueda Semántica, Mapas de Conocimiento.



\newpage
\thispagestyle{empty}
\large{\textbf{Abstract:}}

The research \emph{Recovery, Extraction and Classification of Information from Saber UCV} is presented, where processes of classification, storage and retrieval of information on theses and degree works published in the institutional repository Saber UCV are executed.

A system is implemented that allows classifying the mentioned documents, according to the academic area where the author of the research studied.

The abstracts of the works are extracted and with the obtained classifications a corpus is formed on which an inverted index is generated. Natural Language Processing and Data Mining techniques are applied to the corpus, and with pre-trained artificial intelligence models, text \textit{embeddings} are created. With all the processed information an indexed database is fed.

The system also has a web application for Information Retrieval processes where the user can explore the corpus, including the semantic search, indicating the: text to search, date range, area in which the research was generated, academic level; subsequently the most relevant works are retrieved, enriching the experience with the presentation of the results in interactive tables, Knowledge Maps and recommendations of documents that may be of interest.

The application is implemented under a distributed system with client-server architecture and is supported by the use of orchestrated containers.

\vspace*{2cm}

\textbf{Keywords:} Information Retrieval, Corpus, Natural Language Processing, Relevance,  Artificial Intelligence, Embeddings, Semantic Search, Knowledge Maps.

\thispagestyle{empty}


%\newpage


\setlength{\abovedisplayskip}{-5pt}
\setlength{\abovedisplayshortskip}{-5pt}
\thispagestyle{empty}

\newpage
\begin{center}
\large{\textbf{\emph{\Huge{Dedicatoria:}}}}
\end{center}
\thispagestyle{empty}
\vspace*{5cm}
\thispagestyle{empty}
\begin{center} \Large \emph{A mi hijo Cassiel y  } \end{center}
\vspace*{1cm}
\begin{center} \Large \emph{mi esposa Waleska.} \end{center}
\vspace*{1cm}
\begin{center} \Large {\emph{De ustedes es esta investigación.}} \end{center}



\newpage
\begin{center}
\large{\textbf{\emph{\Huge{Agradecimientos:}}}}
\end{center}
\thispagestyle{empty}
\vspace*{2cm}
\thispagestyle{empty}

- A mi madre, obvio, sino no habría ni una sola palabra acá.\\\\
- A mi padre Fernando por negarme el Atari e insistir en el Oddysey 2.\\\\
- A mi tía Mercedes Infante y mi tío Teófilo García \textdagger, por "lo que fue y será".\\\\
- A mi hermano David por su aguante y soporte.\\\\
- A mi primo Cesar Alejandro García por todo el soporte.\\\\
- Dr. Andres Sanoja primero por aceptar la tutoría y enseñarme qué es la investigación dentro de una comunidad científica.\\\\
- Dr. José Mirabal por siempre andar con alguna idea a desarrollar y el tiempo dedicado a escuchar las propuestas y complementar esta Investigación.\\\\
- Dra. Concettina Di Vasta por las imponentes sesiones de 2 horas 15 minutos llenas de coherencia y conocimiento.\\\\
- Dra. Haydemar Nuñez por la rigurosidad al impartir los conocimientos.\\\\
- Dra. Vanessa Leguizamo por tomarse el tiempo de revisar la solicitud de estudio de un oxidado economista y por ser mi Prof.ª.\\\\
- Dra. Mairene Colina por la ayuda y sugerencias a lo largo de la investigación.\\\\
- Dr. Roberto Abalde por haberme dejado entrar de oyente a su Diplomado y generar tanto interés en estos temas.\\\\
- Mauricio Sáez Toro del equipo Saber UCV por mantener activo el Sistema Saber UCV y disponer del tiempo para colaborar con esta investigación.\\\\
- A Alexandra Asanovna Elbakyan por ayudar a liberar el conocimiento.\\\\
- A todo el personal del Postgrado: sus buenos días, por tener a mano la llave, por ayudar a mantener viva la Academia.\\\\
- A toda la comunidad de creadores de software libre y open source, en especial a los \#useRs por motivarme a adentrarme al mundo de las ciencias de la computación.\\\\


\newpage
\thispagestyle{empty}
\vspace*{5cm}
\hfill
\begin{minipage}{0.70\textwidth}
\begin{quote}
Como todos los hombres de la Biblioteca, he viajado en mi juventud; he peregrinado en busca de un libro, acaso del \emph{catálogo de catálogos}; ahora que mis ojos caso no pueden descifrar lo que escribo, me preparo a morir a unas pocas leguas del hexágono en que nací.\\
--- Jorge Luis Borges, \textit{La Bibioloteca de Babel}, Ficciones
\end{quote}
\hspace*{2cm}

\begin{quote}
Every important aspect of programming arises somewhere in the context of sorting or searching.

--- Donald Knuth, \textit{The Art of Computer Programming}, Volume 3
\end{quote}
\end{minipage}

\thispagestyle{empty}
\maketitle


